\documentclass[10pt]{extarticle}
\usepackage[a4paper,margin=1.5cm]{geometry}
\usepackage[T1]{fontenc}
\usepackage{amsmath,amssymb}
\usepackage{graphicx}
\usepackage[export]{adjustbox}
\usepackage{multicol}
\usepackage{enumitem}
\usepackage{booktabs,array,tabularx}
\setlength{\columnsep}{0.7cm}
\setlength{\parindent}{0pt}
\setlength{\parskip}{3pt}
\pagestyle{empty}
\makeatletter
% figure registry (figdefs.tex is generated by the extract step)
\newcommand{\qfigdef}[4]{\@namedef{qf@file@#1}{#2}\@namedef{qf@w@#1}{#3}\@namedef{qf@wide@#1}{#4}}
\newcommand{\qf@check}[1]{\@ifundefined{qf@file@#1}{\PackageError{score2}{Unknown figure id #1}{}}{}}
\newcommand{\qf@img}[4][]{\includegraphics[width=#2\dimexpr\@nameuse{qf@w@#4}pt\relax,max width=#3,max height=0.45\textheight,#1]{\@nameuse{qf@file@#4}}}
% Full-width questions: finalize wraps every question that contains a wide figure in \qfullbegin...\qfullend,
% so the WHOLE question (stem, figures, options) is set across the page in reading order, then the two columns
% resume. The question's head - its number, stem and first figure (\qfullsplit marks the end of the head) -
% is kept on one page, so a number is never left at a page foot with its figure overleaf; the rest may break
% across pages normally, so no big gaps are left. Without a \qfullsplit, the whole question is kept together
% if it is under 45% of a page. Never write these yourself.
\newif\ifqf@full
\newif\ifqf@boxed
\newsavebox{\qf@box}
\newcommand{\qf@open}{\global\qf@fulltrue\par\global\qf@boxedtrue
  \setbox\qf@box\vbox\bgroup\hsize\textwidth\linewidth\textwidth}
\newcommand{\qf@emit}[1]{\par\egroup\global\qf@boxedfalse
  \ifdim\dimexpr\ht\qf@box+\dp\qf@box\relax<#1\textheight\noindent\usebox{\qf@box}\par\else\unvbox\qf@box\fi}
\newcommand{\qfullsplit}{\ifqf@boxed\qf@emit{0.85}\fi}
\newcommand{\qf@close}{\ifqf@boxed\qf@emit{0.45}\fi\global\qf@fullfalse}
\newcommand{\qfullbegin}{\end{multicols}\qf@open}
\newcommand{\qfullend}{\qf@close\begin{multicols}{2}}
\newcommand{\qfullstart}{\qf@open}
\newcommand{\qfulldone}{\qf@close}
% \qfig[scale]{id}: block figure, centred, directly where it is written (size set by the tool).
% Fallback if a wide figure is not inside a full-width question: break out of the columns in place - no floats.
\newcommand{\qfig}[2][1]{\qf@check{#2}%
  \ifqf@full\qf@blk{#1}{#2}\else
  \ifnum\@nameuse{qf@wide@#2}=1 \expandafter\@firstoftwo\else\expandafter\@secondoftwo\fi
  {\end{multicols}\begin{center}\qf@img{#1}{\textwidth}{#2}\end{center}\begin{multicols}{2}}%
  {\qf@blk{#1}{#2}}\fi}
\newcommand{\qf@blk}[2]{\par\smallskip{\centering\qf@img{#1}{\linewidth}{#2}\par}\smallskip}
% \qfigi[scale]{id}: inline figure (inside option lists / tables / minipages); never breaks columns.
\newcommand{\qfigi}[2][1]{\qf@check{#2}\qf@img[valign=c]{#1}{\linewidth}{#2}}
\makeatother
% structure macros - lint reads these to check completeness and figure placement
\newcommand{\papersec}[1]{\par\bigskip{\centering\Large\bfseries PAPER-#1\par}\medskip}
\newcommand{\partsec}[2]{\par\medskip\noindent{\large\bfseries PART-#1: #2}\par\smallskip}
\newcommand{\sect}[1]{\par\smallskip\noindent\textbf{#1}\par}
\newcommand{\qn}[1]{\par\smallskip\noindent\textbf{#1.}~}
\newcommand{\paraq}[2]{\par\smallskip\noindent\textbf{Paragraph for Questions #1 and #2}\par}
\newenvironment{qoptions}{\begin{itemize}[leftmargin=1.8em,itemsep=1pt,topsep=2pt,label={}]}{\end{itemize}}
\newcommand{\opt}[2]{\item[(#1)] #2}
% figure registry (generated by score2_tools.py finalize)
\qfigdef{p1_phy_s1_q01_1}{TOPS - Test 04 (18 Dec 2022)-img/p1_phy_s1_q01_1.png}{89.9}{0}
\qfigdef{p1_phy_s1_q02_1}{TOPS - Test 04 (18 Dec 2022)-img/p1_phy_s1_q02_1.png}{47.5}{0}
\qfigdef{p1_phy_s1_q03_1}{TOPS - Test 04 (18 Dec 2022)-img/p1_phy_s1_q03_1.png}{205.5}{0}
\qfigdef{p1_phy_s1_q05_1}{TOPS - Test 04 (18 Dec 2022)-img/p1_phy_s1_q05_1.png}{213.9}{0}
\qfigdef{p1_phy_s2_para9_10_1}{TOPS - Test 04 (18 Dec 2022)-img/p1_phy_s2_para9_10_1.png}{107.2}{0}
\qfigdef{p1_phy_s3_para1_2_1}{TOPS - Test 04 (18 Dec 2022)-img/p1_phy_s3_para1_2_1.png}{161.3}{0}
\qfigdef{p1_phy_s3_para5_6_1}{TOPS - Test 04 (18 Dec 2022)-img/p1_phy_s3_para5_6_1.png}{157.0}{0}
\qfigdef{p1_phy_s4_q07_1}{TOPS - Test 04 (18 Dec 2022)-img/p1_phy_s4_q07_1.png}{160.9}{0}
\qfigdef{p1_phy_s4_q08_1}{TOPS - Test 04 (18 Dec 2022)-img/p1_phy_s4_q08_1.png}{104.3}{0}
\qfigdef{p1_phy_s4_q09_1}{TOPS - Test 04 (18 Dec 2022)-img/p1_phy_s4_q09_1.png}{92.3}{0}
\qfigdef{p1_chem_s1_q03_1}{TOPS - Test 04 (18 Dec 2022)-img/p1_chem_s1_q03_1.png}{30.1}{0}
\qfigdef{p1_chem_s1_q03_C}{TOPS - Test 04 (18 Dec 2022)-img/p1_chem_s1_q03_C.png}{24.0}{0}
\qfigdef{p1_chem_s1_q04_1}{TOPS - Test 04 (18 Dec 2022)-img/p1_chem_s1_q04_1.png}{49.9}{0}
\qfigdef{p1_chem_s1_q04_A}{TOPS - Test 04 (18 Dec 2022)-img/p1_chem_s1_q04_A.png}{47.8}{0}
\qfigdef{p1_chem_s1_q04_B}{TOPS - Test 04 (18 Dec 2022)-img/p1_chem_s1_q04_B.png}{47.8}{0}
\qfigdef{p1_chem_s1_q04_C}{TOPS - Test 04 (18 Dec 2022)-img/p1_chem_s1_q04_C.png}{47.8}{0}
\qfigdef{p1_chem_s1_q04_D}{TOPS - Test 04 (18 Dec 2022)-img/p1_chem_s1_q04_D.png}{50.4}{0}
\qfigdef{p1_chem_s2_para9_10_1}{TOPS - Test 04 (18 Dec 2022)-img/p1_chem_s2_para9_10_1.png}{24.0}{0}
\qfigdef{p1_chem_s2_q09_A}{TOPS - Test 04 (18 Dec 2022)-img/p1_chem_s2_q09_A.png}{53.1}{0}
\qfigdef{p1_chem_s2_q09_B}{TOPS - Test 04 (18 Dec 2022)-img/p1_chem_s2_q09_B.png}{42.6}{0}
\qfigdef{p1_chem_s2_q09_C}{TOPS - Test 04 (18 Dec 2022)-img/p1_chem_s2_q09_C.png}{47.8}{0}
\qfigdef{p1_chem_s2_q09_D}{TOPS - Test 04 (18 Dec 2022)-img/p1_chem_s2_q09_D.png}{52.9}{0}
\qfigdef{p1_chem_s2_q10_A}{TOPS - Test 04 (18 Dec 2022)-img/p1_chem_s2_q10_A.png}{47.8}{0}
\qfigdef{p1_chem_s2_q10_B}{TOPS - Test 04 (18 Dec 2022)-img/p1_chem_s2_q10_B.png}{47.2}{0}
\qfigdef{p1_chem_s2_q10_C}{TOPS - Test 04 (18 Dec 2022)-img/p1_chem_s2_q10_C.png}{156.0}{0}
\qfigdef{p1_chem_s3_q03_1}{TOPS - Test 04 (18 Dec 2022)-img/p1_chem_s3_q03_1.png}{68.6}{0}
\qfigdef{p1_chem_s3_q04_1}{TOPS - Test 04 (18 Dec 2022)-img/p1_chem_s3_q04_1.png}{46.0}{0}
\qfigdef{p1_chem_s4_q09_1}{TOPS - Test 04 (18 Dec 2022)-img/p1_chem_s4_q09_1.png}{28.6}{0}
\qfigdef{p2_phy_s1_q01_1}{TOPS - Test 04 (18 Dec 2022)-img/p2_phy_s1_q01_1.png}{108.3}{0}
\qfigdef{p2_phy_s1_q02_1}{TOPS - Test 04 (18 Dec 2022)-img/p2_phy_s1_q02_1.png}{108.3}{0}
\qfigdef{p2_phy_s1_q03_1}{TOPS - Test 04 (18 Dec 2022)-img/p2_phy_s1_q03_1.png}{154.5}{0}
\qfigdef{p2_phy_s1_q05_1}{TOPS - Test 04 (18 Dec 2022)-img/p2_phy_s1_q05_1.png}{70.5}{0}
\qfigdef{p2_phy_s1_q09_1}{TOPS - Test 04 (18 Dec 2022)-img/p2_phy_s1_q09_1.png}{79.3}{0}
\qfigdef{p2_phy_s1_q11_1}{TOPS - Test 04 (18 Dec 2022)-img/p2_phy_s1_q11_1.png}{182.6}{0}
\qfigdef{p2_phy_s1_q12_1}{TOPS - Test 04 (18 Dec 2022)-img/p2_phy_s1_q12_1.png}{108.3}{0}
\qfigdef{p2_phy_s1_q13_1}{TOPS - Test 04 (18 Dec 2022)-img/p2_phy_s1_q13_1.png}{62.0}{0}
\qfigdef{p2_phy_s1_q14_1}{TOPS - Test 04 (18 Dec 2022)-img/p2_phy_s1_q14_1.png}{45.6}{0}
\qfigdef{p2_phy_s1_q15_1}{TOPS - Test 04 (18 Dec 2022)-img/p2_phy_s1_q15_1.png}{56.6}{0}
\qfigdef{p2_phy_s1_q18_1}{TOPS - Test 04 (18 Dec 2022)-img/p2_phy_s1_q18_1.png}{150.6}{0}
\qfigdef{p2_phy_s1_q19_1}{TOPS - Test 04 (18 Dec 2022)-img/p2_phy_s1_q19_1.png}{167.8}{0}
\qfigdef{p2_phy_s1_q20_1}{TOPS - Test 04 (18 Dec 2022)-img/p2_phy_s1_q20_1.png}{84.4}{0}
\qfigdef{p2_phy_s2_q05_1}{TOPS - Test 04 (18 Dec 2022)-img/p2_phy_s2_q05_1.png}{110.4}{0}
\qfigdef{p2_phy_s2_q09_1}{TOPS - Test 04 (18 Dec 2022)-img/p2_phy_s2_q09_1.png}{111.1}{0}
\qfigdef{p2_phy_s2_q10_1}{TOPS - Test 04 (18 Dec 2022)-img/p2_phy_s2_q10_1.png}{76.5}{0}
\qfigdef{p2_chem_s1_q01_1}{TOPS - Test 04 (18 Dec 2022)-img/p2_chem_s1_q01_1.png}{133.7}{0}
\qfigdef{p2_chem_s1_q05_1}{TOPS - Test 04 (18 Dec 2022)-img/p2_chem_s1_q05_1.png}{82.0}{0}
\qfigdef{p2_chem_s1_q07_A}{TOPS - Test 04 (18 Dec 2022)-img/p2_chem_s1_q07_A.png}{138.6}{0}
\qfigdef{p2_chem_s1_q07_B}{TOPS - Test 04 (18 Dec 2022)-img/p2_chem_s1_q07_B.png}{139.4}{0}
\qfigdef{p2_chem_s1_q07_C}{TOPS - Test 04 (18 Dec 2022)-img/p2_chem_s1_q07_C.png}{139.3}{0}
\qfigdef{p2_chem_s1_q07_D}{TOPS - Test 04 (18 Dec 2022)-img/p2_chem_s1_q07_D.png}{138.1}{0}
\qfigdef{p2_chem_s1_q08_1}{TOPS - Test 04 (18 Dec 2022)-img/p2_chem_s1_q08_1.png}{24.0}{0}
\qfigdef{p2_chem_s1_q09_1}{TOPS - Test 04 (18 Dec 2022)-img/p2_chem_s1_q09_1.png}{69.0}{0}
\qfigdef{p2_chem_s1_q10_A}{TOPS - Test 04 (18 Dec 2022)-img/p2_chem_s1_q10_A.png}{32.0}{0}
\qfigdef{p2_chem_s1_q10_C}{TOPS - Test 04 (18 Dec 2022)-img/p2_chem_s1_q10_C.png}{55.3}{0}
\qfigdef{p2_chem_s1_q10_D}{TOPS - Test 04 (18 Dec 2022)-img/p2_chem_s1_q10_D.png}{44.8}{0}
\qfigdef{p2_chem_s1_q11_D}{TOPS - Test 04 (18 Dec 2022)-img/p2_chem_s1_q11_D.png}{138.6}{0}
\qfigdef{p2_chem_s1_q12_1}{TOPS - Test 04 (18 Dec 2022)-img/p2_chem_s1_q12_1.png}{89.9}{0}
\qfigdef{p2_chem_s1_q13_1}{TOPS - Test 04 (18 Dec 2022)-img/p2_chem_s1_q13_1.png}{84.5}{0}
\qfigdef{p2_chem_s1_q13_B}{TOPS - Test 04 (18 Dec 2022)-img/p2_chem_s1_q13_B.png}{58.3}{0}
\qfigdef{p2_chem_s1_q20_1}{TOPS - Test 04 (18 Dec 2022)-img/p2_chem_s1_q20_1.png}{154.2}{0}
\qfigdef{p2_chem_s2_q02_1}{TOPS - Test 04 (18 Dec 2022)-img/p2_chem_s2_q02_1.png}{148.0}{0}
\qfigdef{p2_chem_s2_q03_1}{TOPS - Test 04 (18 Dec 2022)-img/p2_chem_s2_q03_1.png}{129.6}{0}
\qfigdef{p2_chem_s2_q05_0}{TOPS - Test 04 (18 Dec 2022)-img/p2_chem_s2_q05_0.png}{28.2}{0}
\qfigdef{p2_chem_s2_q05_1}{TOPS - Test 04 (18 Dec 2022)-img/p2_chem_s2_q05_1.png}{85.8}{0}
\qfigdef{p2_chem_s2_q05_2}{TOPS - Test 04 (18 Dec 2022)-img/p2_chem_s2_q05_2.png}{139.9}{0}
\qfigdef{p2_chem_s2_q05_3}{TOPS - Test 04 (18 Dec 2022)-img/p2_chem_s2_q05_3.png}{107.3}{0}
\qfigdef{p2_chem_s2_q06_1}{TOPS - Test 04 (18 Dec 2022)-img/p2_chem_s2_q06_1.png}{87.2}{0}
\qfigdef{p2_chem_s2_q07_1}{TOPS - Test 04 (18 Dec 2022)-img/p2_chem_s2_q07_1.png}{44.8}{0}
\qfigdef{p2_chem_s2_q08_1}{TOPS - Test 04 (18 Dec 2022)-img/p2_chem_s2_q08_1.png}{99.2}{0}
\qfigdef{p2_chem_s2_q08_2}{TOPS - Test 04 (18 Dec 2022)-img/p2_chem_s2_q08_2.png}{153.5}{0}
\raggedcolumns
\makeatletter\@beginparpenalty=10000\makeatother
\newcolumntype{L}[1]{>{\raggedright\arraybackslash}p{#1}}
\begin{document}
\begin{center}{\Large\bfseries TOPS --- Test 04 (18 Dec 2022)}\end{center}
\begin{multicols}{2}
\raggedcolumns
\papersec{1}
\partsec{1}{PHYSICS}
\sect{SECTION-I(i)}
\qn{1} A particle of mass $m$ is fixed on a massless ring which rests on a sufficiently rough surface so as to prevent slipping. The system is released from the shown configuration. Then :
\qfig{p1_phy_s1_q01_1}
\begin{qoptions}
\opt{A}{Immediately after releasing from the shown configuration the angular acceleration of the ring is $\dfrac{g\sin\theta}{4R\sin^2\left(\frac\theta2\right)}$.}
\opt{B}{Immediately after releasing from the shown configuration is the angular acceleration of the ring is $\dfrac{g\sin\theta}{2R}$.}
\opt{C}{If $\alpha$ be the anticlockwise angular acceleration immediately after release from the shown position, then the normal reaction immediately after release will be $mg-mR\alpha\cos\theta$}
\opt{D}{If $\alpha$ be the anticlockwise angular acceleration immediately after release from the shown position, then the normal reaction immediately after release will be $mg-mR\alpha\sin\theta$}
\end{qoptions}
\qn{2} A rod is free to rotate in horizontal plane about one of its end hinged at point `$O$' as shown in figure. Point $P$ \& $Q$ are fixed points located in horizontal plane in such way that $PO$ is perpendicular to rod \& $QO$ is along length of rod. A bullet hits the rod perpendicularly at some distance from point `$O$' and gets imbedded into it. Then :
\qfig{p1_phy_s1_q02_1}
\begin{qoptions}
\opt{A}{Linear momentum of rod + bullet system will decrease after collision.}
\opt{B}{Linear momentum of rod + bullet system may increase}
\opt{C}{Angular momentum of rod + bullet system about point `$P$' will not change during the collision.}
\opt{D}{Angular momentum of rod + bullet system about `$Q$' may increase after collision.}
\end{qoptions}
\qn{3} A pendulum bob has mass $m$. The length of pendulum is $\ell$. It is initially at rest. A particle $P$ of mass $m/2$ moving horizontally along $-$ve $x$-direction with velocity $\sqrt{2g\ell}$ collides with the bob and comes to rest. When the bob comes to rest momentarily, another particle $Q$ of mass $m$ moving horizontally along $z$-direction collides with the bob and sticks to it. It is observed that the bob now moves along a horizontal circle.

The floor is a horizontal surface at a distance $2\ell$ below the point of suspension of the pendulum.
\qfig{p1_phy_s1_q03_1}
\begin{qoptions}
\opt{A}{Tension in string immediately after the first collision is $\dfrac32mg$.}
\opt{B}{The height of circular path above the floor is $\dfrac{5\ell}4$}
\opt{C}{Time period of the circular motion $2\pi\sqrt{\dfrac\ell g}$}
\opt{D}{Time period of the circular motion $2\pi\sqrt{\dfrac{3\ell}{4g}}$}
\end{qoptions}
\qn{4} A particle of mass `$m$' starts from origin and move in $x$-$y$ plane. Velocity of particle is given by $\vec v=k_1\hat i+k_2\sqrt{1-\alpha y}\,\hat j$

where $k_1$, $k_2$ \& $\alpha$ are constant and `$y$' is $y$-coordinate of position of particle. Choose the correct statement(s) :
\begin{qoptions}
\opt{A}{Average force acting on particle $\dfrac{m\alpha k_2^2}2$}
\opt{B}{Force acting on particle decreases first then increases}
\opt{C}{Path of particle will be hyperbolic.}
\opt{D}{Radius of curvature of the path followed by particle will decrease first then increase.}
\end{qoptions}
\qn{5} A block of mass 10 kg is put gently on a belt-conveyor system of infinite length at $t=0$ sec, which is moving with constant speed 20 m/sec rightward at all time, irrespectively of any situation by means of a motor-system as shown in the figure. Constant force of magnitude 15 N is applied on the block continuously during its motion. Mark the INCORRECT statement(s) :
\qfig{p1_phy_s1_q05_1}
\begin{qoptions}
\opt{A}{The nature of friction force acting on block of mass 10 kg during its motion is always kinetic in nature.}
\opt{B}{Work done by the kinetic friction on the block of mass 10 kg is 1250 Joule.}
\opt{C}{The magnitude of acceleration of the block of mass 10 kg at $t=6$ sec is 4 m/s$^2$.}
\opt{D}{The magnitude of friction on the block of mass 10 kg at $t=6$ sec is zero.}
\end{qoptions}
\qn{6} A uniform metallic circular disc of mass $M$ and radius $R$, mounted on frictionless bearings, is rotating an angular frequency $\omega$ about an axis passing through its centre and perpendicular to its plane. The temperature of the disc is then increased by $\Delta t$. If $\alpha$ is the coefficient of linear expansion of the metal,
\begin{qoptions}
\opt{A}{The moment of inertia increases by $MR^2\alpha\Delta t$.}
\opt{B}{The moment of inertia remains unchanged.}
\opt{C}{The angular frequency increases by $2\alpha\omega\Delta t$.}
\opt{D}{The angular frequency decreases by $2\alpha\omega\Delta t$.}
\end{qoptions}
\sect{SECTION-I(ii)}
\paraq{7}{8} A uniform rod of mass $m$ and length $\ell$ pivoted at one of its top end is hinging freely in vertical plane. Another identical rod moving horizontally with velocity $v$ along a line passing through its lower end hits it and sticks to it. The two rods were perpendicular during the hit and later also they remain perpendicularly connected to each other. (Use : $v^2=\dfrac{5g\ell}3$)
\qn{7} Find the maximum angle turned by the two-rod system after collision.
\begin{qoptions}
\opt{A}{$\cos^{-1}\left(\dfrac1{\sqrt{10}}\right)+\cos^{-1}\left(\dfrac2{\sqrt{10}}\right)$}
\opt{B}{$\sin^{-1}\left(\dfrac1{\sqrt{10}}\right)+\cos^{-1}\left(\dfrac3{\sqrt{10}}\right)$}
\opt{C}{$\cos^{-1}\left(\dfrac2{\sqrt{10}}\right)+\cos^{-1}\left(\dfrac3{\sqrt{10}}\right)$}
\opt{D}{$\cos^{-1}\left(\dfrac1{\sqrt{10}}\right)+\cos^{-1}\left(\dfrac3{\sqrt{10}}\right)$}
\end{qoptions}
\qn{8} Mark the correct statement :
\begin{qoptions}
\opt{A}{Linear momentum is conserved during collision.}
\opt{B}{After collision speed of combined centre of mass of system continuously decreases till it stops at extreme.}
\opt{C}{After collision speed of combined centre of mass of system first increases then decreases till it stops at extreme.}
\opt{D}{After collision speed of combined centre of mass of system first decreases then increases till it stops at extreme.}
\end{qoptions}
\paraq{9}{10} A light rod of length $L$ is hinged to a plank of mass $m$. The plank is lying on the edge of a horizontal table such that the rod can swing freely in the vertical plane without any hindrance from the table. A particle of mass $m$ is attached to the end of the rod and system is released from $\theta=0^\circ$ position (see figure)
\qfig{p1_phy_s2_para9_10_1}
\qn{9} Assume that friction is absent everywhere and calculate the speed of the plank when the rod makes $\theta=180^\circ$.
\begin{qoptions}
\opt{A}{$\sqrt{8gL}$}
\opt{B}{$\sqrt{6gL}$}
\opt{C}{$2\sqrt{gL}$}
\opt{D}{$\sqrt{2gL}$}
\end{qoptions}
\qn{10} Assume that friction between the plank and the table is large enough to prevent it from slipping and calculate the smallest normal force applied by the plank on the table.
\begin{qoptions}
\opt{A}{$\dfrac{2mg}3$}
\opt{B}{$\dfrac{mg}3$}
\opt{C}{$mg$}
\opt{D}{$\dfrac{mg}2$}
\end{qoptions}
\sect{SECTION-II(i)}
\paraq{1}{2} A disc $D$ and a fix vertical cylinder $C$ of radius $R=0.25$ m connected through a horizontal thread $CD$ of length $\ell=2$m lies in a horizontal plane. On imparting an initial velocity of $v=2$ m/s to the disc it starts moving on the path shown in the figure.
\qfig{p1_phy_s3_para1_2_1}
\qn{1} What time (in sec) will the disc take to strike the cylinder ?
\qn{2} What is the total distance (in meter) travelled by disc before striking the cylinder ?
\paraq{3}{4} A small block of mass $m$ slides down an inclined plane forming an angle $\alpha=30^\circ$ with the horizontal. The coefficient of friction $k$ between the block and the plane is no longer constant but varies with $x$, the distance travelled by the block on the plane, according to the equation $k=ax$, where $a$ is a constant. (Use : $a=\dfrac{\sqrt3}2$, $g=10$ m/s$^2$)
\qn{3} Find the distance (in meter) travelled by the block before coming to rest.
\qn{4} Find the maximum speed (in m/s) of block during the motion.
\paraq{5}{6} A weightless rod of length $\ell$ with a small load of mass $m$ at one of its end is held vertical with its lower end hinged on a horizontal surface. The load touches a wedge of mass $M$ in this position. A slight jerk towards right sets the system in motion (see figure), with rod rotating freely in vertical plane about its lower end. There is no friction. (Use : $g\ell=16$ SI units)
\qfig{p1_phy_s3_para5_6_1}
\qn{5} For what mass ratio $\dfrac Mm$ will the rod form an angle $\theta=\pi/3$ with the vertical at the moment the load separates from the wedge?
\qn{6} What is speed (in m/s) of the wedge at that moment?
\sect{SECTION-II(ii)}
\qn{7} A very stiff bar $(AB)$ of negligible mass is suspended horizontally by two vertical rods as shown in figure. Length of the bar is $2.5L$. The steel rod has length $L$ and cross sectional radius of $r$ and the brass rod has length $2L$ and cross sectional radius of $2r$. A vertically downward force $F$ is applied to the bar at a distance $x$ from the steel rod and the bar remains horizontal. Find the value of $x$ if it is given that ratio of Young's modulus of steel and brass is $\dfrac{Y_s}{Y_B}=2$. (Given : $L=1.6$ m)
\qfig{p1_phy_s4_q07_1}
\qn{8} A ball of mass $m$ starts moving from rest from the point $O$ along the track $OCB$ (see figure). Find the reaction $R$ (in Newton) on the ball at $C$ if the curve $OCB$ is defined by equation $y=h\sin\left(\dfrac{\pi x}\ell\right)$ where $h=\dfrac\ell3$. (Use : $mg=\dfrac{18}{9+2\pi^2}$ Newton)
\qfig{p1_phy_s4_q08_1}
\qn{9} A thin rod of length 3m is being rotated in horizontal plane with 10 rad/sec about its one end. The change in length (in cm) of the rod due to rotation is (Given : Young's modulus of elasticity of the wire, $Y=3\times10^8$ N/m$^2$. Density of the rod is $10^4$ kg/m$^3$)
\qfig{p1_phy_s4_q09_1}
\columnbreak
\partsec{2}{CHEMISTRY}
\sect{SECTION-I(i)}
\qn{1} Which of the following statements on critical constants of gasses are correct?
\begin{qoptions}
\opt{A}{Larger the value of $\dfrac{T_c}{P_c}$ of a gas, larger would be the excluded volume.}
\opt{B}{Critical temperature ($T_c$) of gas is greater than its Boyle temperature ($T_b$).}
\opt{C}{At the critical point in the van der Waals gas isotherm $\left(\dfrac{\partial p}{\partial V_m}\right)_{T_c}=0$}
\opt{D}{All are correct}
\end{qoptions}
\qn{2} In the reaction ;

$\mathrm{Zn+HNO_3(dil)\rightarrow Zn(NO_3)_2+N_2O+H_2O}$

Select the correct option - [Atomic mass of Zinc = 65 u]
\begin{qoptions}
\opt{A}{Zn is acting as reducing agent while $\mathrm{HNO_3}$ is acting as oxidising agent}
\opt{B}{Equivalent weight of Zn $=\dfrac{65}2$}
\opt{C}{For 1 mole Zn, 4 moles of $\mathrm{HNO_3}$ are required}
\opt{D}{For 1 mole $\mathrm{N_2O}$ formation, 4 moles of Zn are required}
\end{qoptions}
\qn{3} Correct acidity order
\begin{qoptions}
\opt{A}{$\mathrm{CF_3OH>CCl_3OH}$}
\opt{B}{$\mathrm{CHCl_3>CHF_3}$}
\opt{C}{\qfigi{p1_chem_s1_q03_1} $>$ \qfigi{p1_chem_s1_q03_C}}
\opt{D}{$\mathrm{Cl_3C{-}CO_2H<F_3C{-}CO_2H}$}
\end{qoptions}
\qn{4} In which of the following compound bond length of C--O bond is greater then \qfigi{p1_chem_s1_q04_1}
\begin{qoptions}
\opt{A}{\qfigi{p1_chem_s1_q04_A}}
\opt{B}{\qfigi{p1_chem_s1_q04_B}}
\opt{C}{\qfigi{p1_chem_s1_q04_C}}
\opt{D}{\qfigi{p1_chem_s1_q04_D}}
\end{qoptions}
\qn{5} If $A_0$ atoms of element $X$ have ionization energy $E_1$ where number of ``$s$'' electrons in $X$ is less than 5 and $B_0$ atoms of element $Y$ have ionization energy $E_2$ where number of `$s$' electrons in $Y$ is less than 4. Then.
\begin{qoptions}
\opt{A}{$\dfrac{E_2}{B_0}$ must be greater than $\dfrac{E_1}{A_0}$}
\opt{B}{Size of $Y$ may be greater than $X$ if $Y$ belongs to group 1 \& $X$ belongs to group 16}
\opt{C}{$\dfrac{E_2}{B_0}$ will be greater than $\dfrac{E_1}{A_0}$ if $Y$ belongs to group 18}
\opt{D}{If $Y$ exists in nature as $Y_2$ \& $X$ exists in nature $X_2$ then the compound $X_2Y_2$ will always be planar}
\end{qoptions}
\qn{6} Which of the following set of order is CORRECT for their indicated properties?
\begin{qoptions}
\opt{A}{Bond length of O--F ; $\mathrm{O_2F_2>OF_2}$}
\opt{B}{Lewis acidic strength ; $\mathrm{BCl_3>BF_3}$}
\opt{C}{Basic strength ; $\mathrm{O(CH_3)_2<O(SiH_3)_2}$}
\opt{D}{$p\pi$--$p\pi$ back bond strength ; $\mathrm{BF_3>BCl_3}$}
\end{qoptions}
\sect{SECTION-I(ii)}
\paraq{7}{8} Potassium dichromate ($\mathrm{K_2Cr_2O_7}$) is an orange coloured compound, very frequently used in laboratory as an oxidising agent as well as in a redox titration. It is generally prepared from chromite ($\mathrm{FeCr_2O_4}$) ore according to the following reactions:

(1) Fusion of chromite ore with sodium carbonate in excess of air.

$\mathrm{FeCr_2O_4+Na_2CO_3+O_2\longrightarrow Na_2CrO_4+Fe_2O_3+CO_2}$

(2) Acidifying filtered sodium chromate solution with sulphuric acid.

$\mathrm{Na_2CrO_4+H_2SO_4\longrightarrow Na_2Cr_2O_7+Na_2SO_4+H_2O}$

(3) Treating sodium dichromate with potassium chloride.

$\mathrm{Na_2Cr_2O_7+KCl\longrightarrow K_2Cr_2O_7+NaCl}$

Answer the following questions using above information. [Fe = 56, Cr = 52]
\qn{7} If you are initially provided with 224 gm of pure chromite ore \& 169.6 gm of sodium carbonate, the minimum volume of air required at 1 atm and 273 K to consume at least one of the reactant completely, if air contains 20\% by volume of oxygen gas is
\begin{qoptions}
\opt{A}{156.8 L}
\opt{B}{31.36 L}
\opt{C}{28 L}
\opt{D}{152.4 L}
\end{qoptions}
\qn{8} If the number of moles of reactants available for reaction are: \{$\mathrm{FeCr_2O_4}$ = 0.25 moles, $\mathrm{O_2}$ = 0.35 moles, $\mathrm{Na_2CO_3}$ = 0.60 moles, $\mathrm{H_2SO_4}$ = 0.2 moles, KCl = 0.1 moles\}, then the maximum number of moles of $\mathrm{K_2Cr_2O_7}$ that can be produced is
\begin{qoptions}
\opt{A}{0.05 moles}
\opt{B}{0.1 moles}
\opt{C}{0.2 moles}
\opt{D}{0.5 moles}
\end{qoptions}
\paraq{9}{10} Planar monocyclic ring containing $(4n+2)p$ electron (where $n=0,1,2,3\ldots$) are aromatic on applying this rule benzene and its derivative are aromatic.

Non benzenoid aromatic compound are those which contain ring that is not six membered such as \qfigi{p1_chem_s2_para9_10_1}.

Cyclic planar conjugated compound which posses only $4np$ electron are antiaromatic in a reaction that reactant would tend to react with reagent that can result into aromatic product salt.
\qn{9} Which of the following is correct option regarding direction of dipole moment ?
\begin{qoptions}
\opt{A}{\qfigi{p1_chem_s2_q09_A}}
\opt{B}{\qfigi{p1_chem_s2_q09_B}}
\opt{C}{\qfigi{p1_chem_s2_q09_C}}
\opt{D}{\qfigi{p1_chem_s2_q09_D}}
\end{qoptions}
\qn{10} Which of the following is correct ?
\begin{qoptions}
\opt{A}{\qfigi{p1_chem_s2_q10_A} reacts immediately with $\mathrm{AgNO_3}$}
\opt{B}{\qfigi{p1_chem_s2_q10_B} reacts immediately with $\mathrm{AgNO_3}$}
\opt{C}{\qfigi{p1_chem_s2_q10_C}}
\opt{D}{Both (A) and (B)}
\end{qoptions}
\sect{SECTION-II(i)}
\paraq{1}{2} Gas `A' is taken in a closed container at the initial total pressure of 1000 torr. Pressure of gas decreases to 900 torr in 5 seconds due to the effusion through a square cross-section. Another similar sized $2^{\text{nd}}$ container is taken in which gaseous mixture of $A$ and $B$ is taken at 5000 torr. Initial molar mass of mixture is $\dfrac{472}5$ gm/mol. A rectangular cross section is made in this container and gases are allowed to effuse. Width of this cross-section is same as the side of the previous square cross-section and length of the rectangular cross-section is 50\% more than that of its width.

(i) Gases $A$ and $B$ are non-reacting

(ii) Rate of effusion is directly proportional to area of orifice.

(iii) Rate of effusion is independent of the change in pressure due to diffusion.

Assume all other conditions to be identical. Answer the following question.

(Molar mass of $A=128$ gm/mole, Molar mass of $B=72$ gm/mole)
\qn{1} mol \% of gas $A$ in mixture effused out initially from $2^{\text{nd}}$ container is :
\qn{2} Change is pressure (in torr) of gas $B$ in container in first 5 seconds due to effusion is-
\paraq{3}{4} Compounds having same molecular formula same structure but different 3D arrangement of atoms or groups are called stereoisomer of each other and phenomenon is called stereo isomerism.
\qn{3} Possible number of G.I. of compound?
\qfig{p1_chem_s3_q03_1}
\qn{4} How many chiral carbon are present in
\qfig{p1_chem_s3_q04_1}
\paraq{5}{6} Element ``$A$'' can form neutral oxides as well as acidic oxide element ``$A$'' forms a polar compound with element ``$B$'' ($Z_{\text{eff}}=Z$) (all possible combinations)
\qn{5} The difference in degeneracy of $2^{\text{nd}}$ excited state of $B$ \& number of $\pi$ bonds in $\mathrm{A_2O}$ is \underline{\hspace{1cm}}.
\qn{6} Sum of sigma bonds in $\mathrm{AB_3}$, $\mathrm{BA_3}$ and $\mathrm{A_2B_4}$ is \underline{\hspace{1cm}}.
\sect{SECTION-II(ii)}
\qn{7} Read the following statements :

(1) If a gas expands at constant temperature kinetic energy of the molecule is decreases.

(2) If a gas expands at constant pressure by providing heat, the kinetic energy of gaseous molecules remains same.

(3) At very low pressure real gases show minimum deviation from ideal behaviour.

(4) The compressibility factor for an ideal gas is zero.

(5) At boyle temperature real gas behave as ideal gas in high pressure region.

(6) Real gas show maximum deviation at high pressure and low temperature.

The number of incorrect statements is :
\qn{8} Number of specie(s) among the following that contain more number of orbitals used by central atom for sigma bond formation as compared to $\mathrm{XeF_2}$ is

$\mathrm{SO_2Cl_2}$, $\mathrm{IO_3^-}$, $\mathrm{H_2O}$, $\mathrm{BeCl_2}$, $\mathrm{NO_2^-}$, $\mathrm{CHCl_3}$, $\mathrm{O_3}$, $\mathrm{NOCl_3}$
\qn{9} Number of species that do NOT exist :-

$\mathrm{NF_5}$, $\mathrm{NOF}$, $\mathrm{PCl_5}$, $\mathrm{OF_6}$, $\mathrm{N_2^-}$, $\mathrm{O_2^{4-}}$, $\mathrm{Be_2}$, $\mathrm{C_2^{2-}}$, $\mathrm{PbI_4}$, $\mathrm{Ne_2}$, $\mathrm{PH_5}$, \qfigi{p1_chem_s4_q09_1}, $\mathrm{NO^+}$
\columnbreak
\partsec{3}{MATHEMATICS}
\sect{SECTION-I(i)}
\qn{1} Let $P(x_1,y_1)$ be a point on the hyperbola $\dfrac{x^2}{a^2}-\dfrac{y^2}{b^2}=1$. The tangent at point $P$ passes through the point $(0,-b)$ and the normal at point $P$ passes through the point $\left(2\sqrt2\,a,\,0\right)$. If $e$ denotes the eccentricity of the hyperbola then which of the following is/are correct ?
\begin{qoptions}
\opt{A}{$x_1=\sqrt2a$}
\opt{B}{$x_1=2a$}
\opt{C}{$y_1=b$}
\opt{D}{$e^2=2$}
\end{qoptions}
\qn{2} Let $(x-k)^2+y^2=r^2$, $(k,r>0)$ is the largest circle, no part of which is lying outside the ellipse $3x^2+4y^2=12$. Now
\begin{qoptions}
\opt{A}{if $k<\dfrac12$ then $r+k=2$}
\opt{B}{if $k>\dfrac12$ then $r+k=2$}
\opt{C}{if $k<\dfrac12$ then $r+k<2$}
\opt{D}{if $k>\dfrac12$ then $r+k<2$}
\end{qoptions}
\qn{3} $x+y=2$ and $x-y=2$ are tangents on a parabola at $(1,1)$ and $(4,2)$ respectively. Which of the following is/are correct
\begin{qoptions}
\opt{A}{Equation of directrix is $x+3y=2$}
\opt{B}{Equation of axis is $3x-y=5$}
\opt{C}{Focus of the parabola is at $\left(\dfrac85,\dfrac65\right)$}
\opt{D}{Vertex of the parabola is at $\left(\dfrac{33}{20},\dfrac{13}{20}\right)$}
\end{qoptions}
\qn{4} Let two parabolas have a common axis where focus of each being exterior to the other and let $\ell_1,\ell_2$ be their latus rectums then the locus of the mid points of the intercepts between the parabolas made on the lines parallel to the common axis is a
\begin{qoptions}
\opt{A}{Straight line if $\ell_1=\ell_2\ \forall\ \ell_1,\ell_2\in R^+$}
\opt{B}{Parabola if $\ell_1\neq\ell_2\ \forall\ \ell_1,\ell_2\in R^+$}
\opt{C}{Parabola $\forall\ \ell_1,\ell_2\in R^+$}
\opt{D}{Hyperbola if $\ell_1=2\ell_2\ \forall\ \ell_1,\ell_2\in R^+$}
\end{qoptions}
\qn{5} If $T_n=\begin{vmatrix}n-1&n-2&n-6\\2n-4&2n-6&2n-11\\3n-9&3n-12&3n-18\end{vmatrix}$ where $n\in N$, then which of the following is/are true ?
\begin{qoptions}
\opt{A}{$\displaystyle\prod_{n=1}^{m}T_n=6^m$}
\opt{B}{$\displaystyle\sum_{n=1}^{10}T_n=60$}
\opt{C}{$\dfrac{T_{n+1}}{T_n}=2$}
\opt{D}{$T_{100}\cdot T_{101}=48$}
\end{qoptions}
\qn{6} If $A(1,3)$ and $C(7,5)$ are two opposite vertices of a square. The equation of a side through $A$ can be -
\begin{qoptions}
\opt{A}{$x+2y=7$}
\opt{B}{$x-2y+5=0$}
\opt{C}{$2x-y+1=0$}
\opt{D}{$2x+y=5$}
\end{qoptions}
\sect{SECTION-I(ii)}
\paraq{7}{8} Consider the inequalities

$I_1:-4\le x^2+y^2-4x-4y\le1$

and $I_2:6x^2-5xy-6y^2-14x+34y-20\ge0$

and draw them in $x$--$y$ plane.
\qn{7} If $(x,y)$ satisfies the inequality $I_1$ then range of values of $x^2+y^2-2x-6y+10$ is equal to
\begin{qoptions}
\opt{A}{$\left[6-4\sqrt2,\ 11+6\sqrt2\right]$}
\opt{B}{$\left[2-\sqrt2,\ 3+\sqrt3\right]$}
\opt{C}{$\left[10-4\sqrt2,\ 10+4\sqrt2\right]$}
\opt{D}{$\left[12-\sqrt2,\ 13+\sqrt2\right]$}
\end{qoptions}
\qn{8} Area common to the inequalities $I_1$ and $I_2$ is
\begin{qoptions}
\opt{A}{$\dfrac52\pi$ sq. units}
\opt{B}{$\dfrac{15}4\pi$ sq. units}
\opt{C}{$\dfrac{13}4\pi$ sq. units}
\opt{D}{$\dfrac94\pi$ sq. units}
\end{qoptions}
\paraq{9}{10} In a triangle $ABC$ if $r_1=2r_2=3r_3$ then (where $r_1,r_2,r_3$ are ex-radius of triangle $ABC$)
\qn{9} The ratio of $\tan\dfrac A2:\tan\dfrac B2:\tan\dfrac C2$ is
\begin{qoptions}
\opt{A}{$6:5:4$}
\opt{B}{$6:3:2$}
\opt{C}{$4:5:6$}
\opt{D}{$3:2:6$}
\end{qoptions}
\qn{10} The ratio of $a:b:c$ is (where $a,b,c$ are sides opposite to angles $A,B,C$ of triangle $ABC$) -
\begin{qoptions}
\opt{A}{$3:4:5$}
\opt{B}{$4:3:5$}
\opt{C}{$5:4:3$}
\opt{D}{$5:3:4$}
\end{qoptions}
\sect{SECTION-II(i)}
\paraq{1}{2} The circle $C_1:x^2+y^2=r^2$ and ellipse $C_2:\dfrac{x^2}{16}+\dfrac{y^2}9=1$ intersect at 4 distinct points $A,B,C,D$. Their common tangents form a parallelogram $A'B'C'D'$.
\qn{1} If $A'B'C'D'$ is a square then $r$ is equal to $\dfrac a{\sqrt b}$ where $a$ and $b$ are coprime then $a+b=$
\qn{2} If $A'B'C'D'$ is a square, then the ratio of area of the circle $C_1$ to the area of the circumcircle of $\Delta A'B'C'$ is $K$, then $K$ is equal to
\paraq{3}{4} Let a straight line $L_1$ whose slope is negative, passing through a point $A(0,8)$, another line perpendicular to $L_1$ meeting it at `$P$' and passing through $B(0,-2)$.
\qn{3} If abscissa of point `$P$' is maximum, then area of the $\Delta APB$ is
\qn{4} If incentre of $\Delta APB$ is $\left(\dfrac{5\sqrt3-5}2,\dfrac{1+5\sqrt3}2\right)$, then equation of the line $L_1$ is $x+\sqrt3y=a\sqrt3$ then value of `$a$' is
\paraq{5}{6} Consider the system of equations $2x+\lambda y+z=0$, $\lambda x+2y+z=0$ and $3x+4y+2z=0$.
\qn{5} For the least value of $\lambda$ if $(x_1,y_1,z_1)$ is a non trivial solution then $\dfrac{x_1+y_1}{z_1}$ is $K$ then value of $14K$ is equal to
\qn{6} The sum of all values of $\lambda$ for which the system has a non trivial solution is $P$ then value of $14P$ is equal to
\sect{SECTION-II(ii)}
\qn{7} Number of value of $x$ in $[0,2\pi]$ for which $2^{(2\sin x-1)(\sin x-1)}+2^{(2\sin x+1)(2-\sin x)}=9$, is
\qn{8} The positive real number which satisfy the equation $\dfrac1{x^2+4x-50}+\dfrac1{x^2+4x-66}=\dfrac2{x^2+4x-90}$ is
\qn{9} Let $A_1,A_2,A_3,\ldots A_{11}$ be 11 arithmetic means and $G_1,G_2,G_3,\ldots G_{11}$ be 11 geometric means and $H_1,H_2,H_3,\ldots H_{11}$ be 11 harmonic means between two positive numbers $a$ and $b$. If $ab=9$ and value of $\displaystyle\prod_{k=1}^{5}\left(A_K\cdot G_{12-2K}\cdot H_{12-K}\right)$ is $p^q$ where $p$ is a prime number and $q$ is a natural number then value of $\left[\dfrac qp\right]$ is (where $[\cdot]$ is greatest integer function).
\columnbreak
\papersec{2}
\partsec{1}{PHYSICS}
\sect{SECTION-I}
\qn{1} Three blocks $A$, $B$ and $C$ of equal mass $m$ are placed one over other on a frictionless surface (table) as shown in the figure. Coefficient of friction between any blocks $A$, $B$ and $C$ is $\mu$. The maximum value of mass of block $M_D$ so that the block $A$, $B$ and $C$ move without slipping over each other is :-
\qfig{p2_phy_s1_q01_1}
\begin{qoptions}
\opt{A}{$\dfrac{3m\mu}{\mu+1}$}
\opt{B}{$\dfrac{3m(1-\mu)}\mu$}
\opt{C}{$\dfrac{3m(1+\mu)}\mu$}
\opt{D}{$\dfrac{3m\mu}{(1-\mu)}$}
\end{qoptions}
\qn{2} A fixed U-shaped smooth wire has a semi-circular bending between $A$ and $B$ as shown in the figure. A bead of mass `$m$' moving with uniform speed $v$ through the wire enters the semicircular bend at $A$ and leaves at $B$. The magnitude of average force exerted by the bead on the part $AB$ of the wire is :-
\qfig{p2_phy_s1_q02_1}
\begin{qoptions}
\opt{A}{0}
\opt{B}{$\dfrac{4mv^2}{\pi d}$}
\opt{C}{$\dfrac{2mv^2}{\pi d}$}
\opt{D}{None of these}
\end{qoptions}
\qn{3} An infinite number of bricks are placed one over the other as shown in the figure. Each succeeding brick having half the length and breadth of its preceding brick and the mass of each succeeding bricks being $1/4^{\text{th}}$ of the preceding one. Take `$O$' as the origin, the $x$-coordinate of centre of mass of the system of bricks is at :-
\qfig{p2_phy_s1_q03_1}
\begin{qoptions}
\opt{A}{$-\dfrac a7$}
\opt{B}{$\dfrac{3a}7$}
\opt{C}{$-\dfrac{3a}7$}
\opt{D}{$-\dfrac{2a}7$}
\end{qoptions}
\qn{4} A particle $m$ starts with zero velocity along a line $y=4d$. The position of particle $m$ varies as $x=A\sin\omega t$. At $\omega t=\pi/2$, its angular momentum with respect to the origin :-
\begin{qoptions}
\opt{A}{$mA\omega d$}
\opt{B}{$m\omega d/A$}
\opt{C}{$mAd/\omega$}
\opt{D}{zero}
\end{qoptions}
\qn{5} A massless string is wrapped around a hollow cylinder having mass $m$ and radius $r$. The cylinder is kept on a rough horizontal surface (coefficient of friction is $\mu$). A constant force $F$ is applied as shown in the figure. In case of pure rolling, the friction force acting on the bottom most point of the cylinder is :-
\qfig{p2_phy_s1_q05_1}
\begin{qoptions}
\opt{A}{zero}
\opt{B}{$\mu mg$}
\opt{C}{$\dfrac{\mu mg}2$}
\opt{D}{$\dfrac{3\mu mg}2$}
\end{qoptions}
\qn{6} If the momentum of a body is increased by 25\% the kinetic energy is increased by :-
\begin{qoptions}
\opt{A}{25\%}
\opt{B}{5\%}
\opt{C}{56\%}
\opt{D}{38\%}
\end{qoptions}
\qn{7} The potential energy of a 1 kg particle free to move along the $x$-axis is given by $V_x=\left(\dfrac{x^4}4-\dfrac{x^2}2\right)$ J. The total mechanical energy is 2J. Then the maximum speed in ms$^{-1}$ is :-
\begin{qoptions}
\opt{A}{2}
\opt{B}{$\dfrac3{\sqrt2}$}
\opt{C}{$\sqrt2$}
\opt{D}{$\dfrac1{\sqrt2}$}
\end{qoptions}
\qn{8} A particle of mass $m$ is acted upon by a force $F=F_0\cos\left(\dfrac\pi2t\right)$. The distance travelled by the particle when it stops first (initial velocity of the particle is zero) is :
\begin{qoptions}
\opt{A}{$\dfrac{2F_0}{m\pi}$}
\opt{B}{$\dfrac{4F_0}{m\pi^2}$}
\opt{C}{$\dfrac{8F_0}{m\pi^2}$}
\opt{D}{zero}
\end{qoptions}
\qn{9} By applying a force $\vec F=(3xy-5z)\hat j+4z\hat k$ a particle is moved along the path $y=x^2$ from point $(0,0,0)$ to the point $(2,4,0)$. The work done by the $F$ on the particle is :-
\qfig{p2_phy_s1_q09_1}
\begin{qoptions}
\opt{A}{$\dfrac{280}5$}
\opt{B}{$\dfrac{140}5$}
\opt{C}{$\dfrac{232}5$}
\opt{D}{$\dfrac{192}5$}
\end{qoptions}
\qn{10} A particle is moving in an isolated $x$--$y$ plane. At an instant, the particle has velocity $(4\hat i+4\hat j)$ m/s and acceleration $(3\hat i+5\hat j)$ m/s$^2$. At that instant what will be the radius of curvature of its path ?
\begin{qoptions}
\opt{A}{$16\sqrt2$ m}
\opt{B}{15 m}
\opt{C}{18 m}
\opt{D}{None of these}
\end{qoptions}
\qn{11} There are two incline planes $AO$ and $OB$ inclined at $45^\circ$ and $60^\circ$ respectively with the horizontal as shown in the figure. When a man moves on the inclined plane $AO$, he observes that the rain drops are falling at an angle $45^\circ$ with the vertical. When the man moves on the inclined plane $OB$, he observes that the raindrops are falling vertically downward with the horizontal Then actual speed of rain is :-
\qfig{p2_phy_s1_q11_1}
\begin{qoptions}
\opt{A}{$5\sqrt2$ m/s}
\opt{B}{$20\sqrt3$ m/s}
\opt{C}{$10\sqrt2$ m/s}
\opt{D}{None of these}
\end{qoptions}
\qn{12} System is shown in the figure. Assume that cylinder remain in contact with the two wedges. The speed of cylinder is :
\qfig{p2_phy_s1_q12_1}
\begin{qoptions}
\opt{A}{$\sqrt{19-4\sqrt3}\,\dfrac u2$ m/s}
\opt{B}{$\dfrac{\sqrt{13}\,u}2$ m/s}
\opt{C}{$\sqrt3u$ m/s}
\opt{D}{$\sqrt7u$ m/s}
\end{qoptions}
\qn{13} A closed cylinder of length `$\ell$' containing a liquid of variable density $\rho(x)=\rho_0(1+\alpha x)$ is rotating with constant angular velocity $\omega$. Find the net force exerted by the liquid on the axis of rotation. (Take the cylinder to be massless and $A$ = cross sectional area of cylinder)
\qfig{p2_phy_s1_q13_1}
\begin{qoptions}
\opt{A}{$\rho_0A\omega^2\ell^2\left[\dfrac12+\dfrac13\alpha\ell\right]$}
\opt{B}{$\rho_0A\omega^2\ell^2\left[\dfrac12+\dfrac23\alpha\ell\right]$}
\opt{C}{$\rho_0A\omega^2\ell^2\left[\dfrac12+\alpha\ell\right]$}
\opt{D}{$\rho_0A\omega^2\ell^2\left[\dfrac12+\dfrac43\alpha\ell\right]$}
\end{qoptions}
\qn{14} A uniform flexible chain of length 1.5 m is placed on a smooth spherical surface of radius $R=\left(\dfrac2\pi\right)$ m with its upper end fixed at the top of the sphere while the other end of the chain is hanging freely as shown in figure. Find the acceleration of each element of the chain when its upper end is released.
\qfig{p2_phy_s1_q14_1}
\begin{qoptions}
\opt{A}{$\left(1+\dfrac4\pi\right)\dfrac g3$}
\opt{B}{$\left(1+\dfrac2\pi\right)\dfrac g6$}
\opt{C}{$\left(1-\dfrac2\pi\right)\dfrac g3$}
\opt{D}{$\left(1+\dfrac3\pi\right)\dfrac g2$}
\end{qoptions}
\qn{15} A cube of side $a$ and mass $M$ is pivoted on a horizontal axis such that it can rotate about one of its edges. Initially it is held in the highest position and then released. Find the net reaction of pivot on the cube at the moment when the reaction towards the center of the cube becomes zero. Given that $MI$ of the cube about its edge is equal to $\dfrac23Ma^2$.
\qfig{p2_phy_s1_q15_1}
\begin{qoptions}
\opt{A}{$\dfrac{Mg}4$}
\opt{B}{$\dfrac{Mg}2$}
\opt{C}{$\dfrac{Mg}5$}
\opt{D}{None of these}
\end{qoptions}
\qn{16} Two boys ``play catch'' with a ball in a long hallway. The ceiling height is $H$, and the ball is thrown and caught at shoulder height, which we call $h$ for each boy. If the boys are capable of throwing the ball with velocity $v_0$, at what maximum separation can they play ?
\begin{qoptions}
\opt{A}{$4\sqrt{(H-h)\left(\dfrac{v_0^2}{2g}+H+h\right)}$}
\opt{B}{$4\sqrt{(H-h)\dfrac{v_0^2}{2g}+(H-h)}$}
\opt{C}{$4\sqrt{(H-h)\left(\dfrac{v_0^2}{2g}-(H+h)\right)}$}
\opt{D}{$4\sqrt{(H-h)\left(\dfrac{v_0^2}{2g}-(H-h)\right)}$}
\end{qoptions}
\qn{17} A cyclist moves along a horizontal circular road of radius $R$, the friction coefficient being dependent only on the distance $r$ from the centre of the circular road as $\mu=\mu_0\left(1-\dfrac rR\right)$, where $\mu_0$ is a constant. What is the value of $r$ for which the cyclist can move with maximum speed without skidding ?
\begin{qoptions}
\opt{A}{$\dfrac R2$}
\opt{B}{$\dfrac R3$}
\opt{C}{$\dfrac R4$}
\opt{D}{$\dfrac R5$}
\end{qoptions}
\qn{18} A uniform thin rod of length $6a$ and mass $8m$ lies on a smooth horizontal table. Two point masses $m$ and $2m$, moving in the same horizontal plane with speeds $2v$ and $v$ respectively, strike the bar as shown in figure and stick to it. The total energy of the system after the collision is :-
\qfig{p2_phy_s1_q18_1}
\begin{qoptions}
\opt{A}{$\dfrac25mv^2$}
\opt{B}{$\dfrac35mv^2$}
\opt{C}{$\dfrac27mv^2$}
\opt{D}{None of these}
\end{qoptions}
\qn{19} A rod $OA$ of mass `$m$' and length $L$ hinged at end $O$ can rotate smoothly in vertical plane about horizontal axis passing through $O$- It is released from a horizontal position as shown in the figure and strikes a block of mass $M$ kept vertically below $O$ at distance $L$ on a smooth horizontal surface. The block after collision moves on the smooth vertical circular track of diameter $L$ such that the force exerted by the circular track on the block at point $C$ is one third of its weight. If the rod after collision reaches position $OP$ when it comes to rest, then find the ratio of masses $M/m$.
\qfig{p2_phy_s1_q19_1}
\begin{qoptions}
\opt{A}{$\dfrac{(\sqrt3+1)}{\sqrt7}$}
\opt{B}{$\dfrac{(\sqrt2+1)}{\sqrt5}$}
\opt{C}{$\dfrac{(\sqrt2+1)}{\sqrt7}$}
\opt{D}{None of these}
\end{qoptions}
\qn{20} A sphere of radius $R$ is projected with a reverse spin $\omega_0$ down a rough inclined plane with a speed $v_0$ for which coefficient of friction is $\mu>\tan\theta$, where $\theta$ is angle of inclination. Minimum value of $\omega_0$ so that it will change the direction of velocity of centre is:-
\qfig{p2_phy_s1_q20_1}
\begin{qoptions}
\opt{A}{$\dfrac{5v_0\mu}{2R(\mu-\tan\theta)}$}
\opt{B}{$\dfrac{3v_0\mu}{5R(\mu-\tan\theta)}$}
\opt{C}{$\dfrac{2v_0\mu}{3R(\mu-\tan\theta)}$}
\opt{D}{$\dfrac{5v_0\mu}{3R(\mu-\tan\theta)}$}
\end{qoptions}
\sect{SECTION-II}
\qn{1} A ladder of mass 20 kg and length 10 m rests against a slippery vertical wall at an angle of $30^\circ$ with the vertical. The ladder, of uniform construction, is prevented from slipping by friction with the ground. What is the magnitude (in N) of the force exerted by the ladder on the wall? ($g=10$ m/s$^2$)
\qn{2} Find the minimum rate (kg/sec) at which the fuel must be consumed by a rocket so as to be able to take off vertically. Given that the total mass at the time of take off is 3000 kg and the velocity of fuel expelled is 300 m/s. [$g=9.8$ m/s$^2$]
\qn{3} A particle describes a horizontal circle on a smooth inner surface of an inverted cone. The height of the plane of the circle above the vertex is 9.8 cm. Find the speed (in m/s) of the particle. (Roundoff the answer to closest integer.) ($g=9.8$ m/s$^2$)
\qn{4} A car starts moving with constant tangential acceleration $a_t=\sqrt2$ m s$^{-2}$ along a horizontal circular road of radius $r=40$ m. The coefficient of sliding friction between the wheels of the car and the road is $\mu=0.20$. Find the distance (in m) moved by the car without sliding. (Round off the answer to closest integer.) ($g=10$ m/s$^2$)
\qn{5} A single string $ACB$ of length 2 m passes through a ring at $C$ that is attached to a sphere which revolves at a constant speed $v$ in the horizontal circle as shown in Figure. Given $\theta_1=60^\circ$, $\theta_2=30^\circ$, $g=10$ ms$^{-2}$, the string is light and inextensible, and tension is same in both parts of the string, If the speed $v$ (in m/s) is $\sqrt{\dfrac{10\sqrt a}{\sqrt b+1}}$. Then find $a^2+b^2$.
\qfig{p2_phy_s2_q05_1}
\qn{6} A sphere of radius 10 cm and mass 25 kg is attached to the lower end of a steel wire of radius 0.05 cm and length 500 cm which is suspended from the ceiling of a room. The point of support is 521 cm above the floor. When the sphere is set swinging as a simple pendulum, its lowest point just grazes the floor, if the velocity of the ball at its lowest position (in m/s) is $v$ then find $v^2$. Given : $Y$ for steel $=20\times10^{10}$ Nm$^{-2}$, $g=9.8$ ms$^{-2}$. Round off the answer to closest integer.
\qn{7} An ice cube at $0^\circ$C of mass 0.1 kg is placed in an isolated container which is at $227^\circ$C. The specific heat capacity $c$ of the material of the container varies with temperature $T$ according to the empirical relation $C=A+BT$, where $A=100$ cal kg$^{-1}$ K$^{-1}$ and $B=2\times10^{-2}$ cal kg$^{-1}$ K$^{-2}$. If final temperature of the container is $27^\circ$C, find the mass of the container (in grams). The specific latent heat of fusion of ice is $8\times10^4$ cal kg$^{-1}$ and the specific heat capacity of water is $10^3$ cal kg$^{-1}$ ($^\circ$C)$^{-1}$.
\qn{8} A particle $P$ is initially at a distance $d=16$ m from a fixed point $O$. The particle $P$ moves with a velocity $\vec v=5\hat{PO}+3\hat i$. Where $\hat{PO}$ is a unit vector from $P$ to $O$ at any time $t$. Initially $\hat{PO}$ is perpendicular to $\hat i$. Find the time in seconds after which point $P$ meets point $O$.
\qn{9} Figure shows a smooth cylindrical pulley of radius $R$ with centre at origin of co-ordinates. An ideal thread is thrown over it on the two parts of ideal thread two identical masses are tied initially at rest with co-ordinates $(R,0)$ and $(-R,-R)$ respectively. If mass at $x$-axis is given a slight upward jerk, it leaves contact with pulley at $(R\cos\phi\ R\sin\phi)$. Then find $\phi/\sin\phi$.
\qfig{p2_phy_s2_q09_1}
\qn{10} An elastic cord having an upstretched length $\ell$, force constant $k$ and mass per unit length $m_0$ is stretched around the drum of radius $R(2\pi R>1)$. Determine the speed of the cord, due to rotation of the drum, which will allow the cord to loosen its contact with the drum. Express value of $\dfrac{v^2}{11}$ (in m$^2$/s$^2$). (for the given data $m_0=40$ g/cm, $k=100$ N/m, $\pi=\dfrac{22}7$, $\ell=400$ cm, $R=70$ cm)
\qfig{p2_phy_s2_q10_1}
\columnbreak
\partsec{2}{CHEMISTRY}
\sect{SECTION-I}
\qn{1} Three containers are containers are connected by narrow tube as shown.
\qfig{p2_chem_s1_q01_1}
Valve I is opened for long period of time and then closed. After it valve II is opened which results in immediate reaction as following :

$\mathrm{NH_3(g)+HCl(g)\longrightarrow NH_4Cl(s)}$

After long time valve II is closed. Calculate the final pressure (in atm) in container `B'.
\begin{qoptions}
\opt{A}{6 atm}
\opt{B}{7 atm}
\opt{C}{8 atm}
\opt{D}{9 atm}
\end{qoptions}
\qn{2} An aqueous solution contains 2M-$\mathrm{AlCl_3}$ and 3m-$\mathrm{SnCl_4}$. If the molarity and molality of $\mathrm{Cl^-}$ ion in the solution is same, the density of solution is (Al = 27, Sn = 119)
\begin{qoptions}
\opt{A}{1.05 gm/ ml}
\opt{B}{1.15 gm/ ml}
\opt{C}{1.95 gm / ml}
\opt{D}{2.05 gm/ ml}
\end{qoptions}
\qn{3} A 150 mL of solution of $\mathrm{I_2}$ is divided into two unequal parts. I part reacts with hypo solution in acidic medium. 15 mL of 0.4 M hypo was consumed. II part was added with 100 mL of 0.3 M NaOH solution. Residual base required 10 mL of 0.3 M $\mathrm{H_3PO_3}$ solution for complete neutralization. What was the initial concentration of $\mathrm{I_2}$ ?
\begin{qoptions}
\opt{A}{0.08 M}
\opt{B}{0.1 M}
\opt{C}{0.2 M}
\opt{D}{none of these}
\end{qoptions}
\qn{4} What is the mass of $\mathrm{Ca(OH)_2}$ required for 10 litre of water to remove temporary hardness of 100 ppm due to $\mathrm{Ca(HCO_3)_2}$ ?
\begin{qoptions}
\opt{A}{1.62 g}
\opt{B}{0.74 g}
\opt{C}{7.4 g}
\opt{D}{0.46 g}
\end{qoptions}
\columnbreak
\qn{5} \qfig{p2_chem_s1_q05_1}
(ii) For liquid $B$.

(iii) For liquid $C$.

Out of liquid $A$, $B$ \& $C$, identify correct option :
\begin{qoptions}
\opt{A}{Liquid $A$, $B$ \& $C$ can be water, ethanol, diethyl ether respectively.}
\opt{B}{Liquid $A$, $B$ \& $C$ can be diethyl ether, ethanol \& water respectively.}
\opt{C}{At temperature $T_1$, liquid $C$ is most volatile.}
\opt{D}{If temperature $>T_2$ \& pressure is 1 atm, then vapour phase of liquid $A$ will not exist.}
\end{qoptions}
\qn{6} A 4g mixture of $\mathrm{P_4}$ and $\mathrm{O_2}$ react completely (No reactant left unreacted) and form equal moles of $\mathrm{P_4O_6}$ and $\mathrm{P_4O_{10}}$. Find mass\% of $\mathrm{P_4}$ in original sample. (P = 31)
\begin{qoptions}
\opt{A}{60.2\%}
\opt{B}{51.4\%}
\opt{C}{49.2\%}
\opt{D}{30.4\%}
\end{qoptions}
\qn{7} Which of the following is CORRECT :
\begin{qoptions}
\opt{A}{\qfigi{p2_chem_s1_q07_A} hydrogenation}
\opt{B}{\qfigi{p2_chem_s1_q07_B}}
\opt{C}{\qfigi{p2_chem_s1_q07_C}}
\opt{D}{\qfigi{p2_chem_s1_q07_D} Heat of hydrogenation}
\end{qoptions}
\qn{8} Which compound is most basic -
\begin{qoptions}
\opt{A}{$\mathrm{CH_3{-}CH_2{-}NH_2}$}
\opt{B}{$\mathrm{CH_3{-}CH{=}NH}$}
\opt{C}{$\mathrm{CH_3{-}C{\equiv}N}$}
\opt{D}{\qfigi{p2_chem_s1_q08_1}}
\end{qoptions}
\qn{9} Arrange marked hydrogen in decresing order of their acidic strength in vitamin C.
\qfig{p2_chem_s1_q09_1}
\begin{qoptions}
\opt{A}{I $>$ II $>$ III}
\opt{B}{III $>$ II $>$ I}
\opt{C}{II $>$ III $>$ I}
\opt{D}{II $>$ I $>$ III}
\end{qoptions}
\qn{10} Which of the following compound will show tautomerism?
\begin{qoptions}
\opt{A}{\qfigi{p2_chem_s1_q10_A}}
\opt{B}{$\mathrm{Me_3C{-}C{\equiv}N}$}
\opt{C}{\qfigi{p2_chem_s1_q10_C}}
\opt{D}{\qfigi{p2_chem_s1_q10_D}}
\end{qoptions}
\qn{11} Select the correct statement -
\begin{qoptions}
\opt{A}{Homologoues can not be represented by same general formula}
\opt{B}{$\mathrm{C_5H_{11}}$--- has seven alkyl groups}
\opt{C}{$\mathrm{C_4H_9}$--- has four alkyl groups}
\opt{D}{\qfigi{p2_chem_s1_q11_D}}
\end{qoptions}
\qn{12} Select the correct name for the following diene
\qfig{p2_chem_s1_q12_1}
\begin{qoptions}
\opt{A}{2-(E)--4--(Z)--octa--2,4--diene}
\opt{B}{2-(E)--4--(E)--octa--2,4--diene}
\opt{C}{2-(Z)--4--(E)--octa--2,4--diene}
\opt{D}{2-(Z)--4--(Z)--octa--2,4--diene}
\end{qoptions}
\qn{13} (A) \qfig{p2_chem_s1_q13_1}
(B) \qfig{p2_chem_s1_q13_B}
Find out the correct relation in the above pair of compound?
\begin{qoptions}
\opt{A}{Diastereomers}
\opt{B}{Enantiomers}
\opt{C}{Position isomers}
\opt{D}{Conformation}
\end{qoptions}
\qn{14} The correct mentioned order against the indicated property is.
\begin{qoptions}
\opt{A}{$\mathrm{SbH_3>NH_3>PH_3>AsH_3}$ (Boiling point)}
\opt{B}{$\mathrm{H_2O>H_2Te>H_2Se>H_2S}$ (Boiling point)}
\opt{C}{o-hydroxyphenol $>$ pyrogallol (Dipole moment)}
\opt{D}{$\mathrm{P_4O_6<P_4O_{10}}$ (X--O--X linkages)}
\end{qoptions}
\qn{15} Ionic mobility of an ion in an aqueous solution is given as

$\text{Ionic Mobility}=\dfrac{\text{Speed of an ion}}{\text{Potential gradient}}$

Assuming ions such as $\mathrm{Na^+}$, $\mathrm{Al^{3+}}$, $\mathrm{K^+}$ and $\mathrm{Ti^{4+}}$ are moving in aqueous solutions through same potential gradient the time required to reach the respective electrode from same starting point will be maximum for
\begin{qoptions}
\opt{A}{$\mathrm{Ti^{4+}}$}
\opt{B}{$\mathrm{K^+}$}
\opt{C}{$\mathrm{Na^+}$}
\opt{D}{$\mathrm{Al^{3+}}$}
\end{qoptions}
\qn{16} \par\smallskip
{\footnotesize\setlength{\tabcolsep}{3pt}
\begin{tabular}{|l|p{2.5cm}|l|p{2.6cm}|}
\hline
\multicolumn{2}{|c|}{List - I (Species)} & \multicolumn{2}{c|}{List - II (Type of back bonding)}\\
\hline
(P) & $\mathrm{(CH_3)_3CO^-}$ & (1) & $np_\pi-np_\pi$\\
(Q) & Boron trichloride & (2) & $np_\pi-(n+1)p_\pi$\\
(R) & Phosphorous trifluoride & (3) & No back bonding\\
(S) & Diflorocarbene (singlate) & (4) & $np_\pi-(n+1)d_\pi$\\
\hline
\end{tabular}}
\begin{qoptions}
\opt{A}{P$\to$3 ; Q$\to$2 ; R$\to$1 ; S$\to$4}
\opt{B}{P$\to$1 ; Q$\to$2 ; R$\to$4 ; S$\to$3}
\opt{C}{P$\to$1 ; Q$\to$4 ; R$\to$2 ; S$\to$3}
\opt{D}{P$\to$3 ; Q$\to$2 ; R$\to$4 ; S$\to$1}
\end{qoptions}
\qn{17} Which of the following order is correct?
\begin{qoptions}
\opt{A}{AgCl $<$ AgBr (Solubility in water)}
\opt{B}{$\mathrm{BeCO_3>BaCO_3}$ (Thermal Stability)}
\opt{C}{$\mathrm{K^+<Cu^+}$ (Polarizing power)}
\opt{D}{KF $>$ NaF (Covalent character)}
\end{qoptions}
\qn{18} The INCORRECT order of boiling point is
\begin{qoptions}
\opt{A}{$\mathrm{H_2O>CH_3OH}$}
\opt{B}{$\mathrm{N(CH_3)H_2<CH_3OH}$}
\opt{C}{$\mathrm{H_3PO_4>Me_3PO_4}$}
\opt{D}{$\mathrm{CH_3N_3>HN_3}$}
\end{qoptions}
\qn{19} In $\mathrm{NO_2}$ molecule :-
\begin{qoptions}
\opt{A}{Bond angle is lower than $120^\circ$}
\opt{B}{Diamagnetic nature}
\opt{C}{Odd number of $e^-$ in valence shell}
\opt{D}{both N--O length are identical}
\end{qoptions}
\qn{20} Consider the given L.C.A.O. \& find the correct statement :-
\qfig{p2_chem_s1_q20_1}
\begin{qoptions}
\opt{A}{Resulting M.O. is $\sigma$ antibonding}
\opt{B}{Resulting M.O. is ungerade}
\opt{C}{Has no nodal plane which is perpendicular to the inter-nuclear axis \& passing between the nucleus of both}
\opt{D}{resulting M.O. is H.O.M.O. of $\mathrm{B_2}$}
\end{qoptions}
\sect{SECTION-II}
\qn{1} Beautician of the Queen Elizabeth had an $\mathrm{H_2O_2}$ solution of molality `5m' \& density 1.17 gm/ml. If 2.5 mol of $\mathrm{O_2}$ is needed for bleaching of the Queen's hair, then find the minimum mass (in grams) of the $\mathrm{H_2O_2}$ solution the beautician needs for the bleaching purpose. Do not forget that the $\mathrm{H_2O_2}$ decomposes to give water and oxygen gas.
\qn{2} The graph of compressibility factor $(Z)$ vs. $P$ for one mole of a real gas is shown in following diagram. The graph is plotted at constant temperature 273 K. If the slope of graph at very high pressure $\left(\dfrac{dZ}{dP}\right)$ is $\left(\dfrac1{2.8}\right)$ atm$^{-1}$, then calculate volume of one mole of real gas molecules (in L / mol)

Given : $N_A=6\times10^{23}$ and $R=\dfrac{22.4}{273}$ L atm K$^{-1}$mol$^{-1}$
\qfig{p2_chem_s2_q02_1}
\qn{3} Two moles of an ideal gas undergoes the following process. Given that $\left(\dfrac{\partial P}{\partial T}\right)_V$ is $x\times10^{-y}$ atm/k, then calculate the value of $(x+y)$
\qfig{p2_chem_s2_q03_1}
\qn{4} A 10 gm sample of NaCN is dissolved in 50 ml of 0.3 M mild alkaline solution of $\mathrm{KMnO_4}$ and heated strongly to convert all $\mathrm{CN^-}$ to $\mathrm{OCN^-}$. The left solution is then acidified with $\mathrm{H_2SO_4}$ which require 500 ml of 0.05 M $\mathrm{FeSO_4}$. Calculate percentage purity of NaCN in sample.
\qn{5} (a) \qfigi{p2_chem_s2_q05_0}\hspace{0.5em}\qfigi{p2_chem_s2_q05_1}
\qfig{p2_chem_s2_q05_2}
\qfig{p2_chem_s2_q05_3}
For the above compound identify the number

$a=x-y$

$x=$ Number of compounds which release $\mathrm{CO_2}$ on reaction with $\mathrm{NaHCO_3}$

$y=$ Number of compounds which are aromatic
\qn{6} How many chiral centres are present in the following compound ?
\qfig{p2_chem_s2_q06_1}
\qn{7} How many geometrical isomerse are possible in the given compound?
\qfig{p2_chem_s2_q07_1}
\qn{8} Consider the following

(a) $\mathrm{N_2>O_2}$ \quad $1^{\text{st}}$ ionisation enthalpy

(b) $\mathrm{N_2>N}$ \quad $1^{\text{st}}$ ionisation enthalpy

(c) $\mathrm{O_2<N}$ \quad $1^{\text{st}}$ ionisation enthalpy

(d) $\mathrm{O_2<O}$ \quad $1^{\text{st}}$ ionisation enthalpy

(e) $\mathrm{O_2>N_2}$ \quad Bond order

(f) $\mathrm{N_2>O_2}$ \quad Bond length

(g) $\mathrm{N_2=O_2}$ \quad Bond strength
\qfig{p2_chem_s2_q08_1}
After LCAO there will be one new nodal plane formation take place in resulting molecular orbital.

(i) $\mathrm{CO^+}$ and $\mathrm{N_2^{2+}}$ have same bond order.
\qfig{p2_chem_s2_q08_2}
gerade orbital is produced

How many of the above statements are correct:
\qn{9} The spin only magnetic moment of an atom of element ``$A$'' in its ground state is 1.73 BM. This element belongs to p-block and has significantly high EA value. Another element ``$B$'' of same group and having less EA than $A$ forms a compound $\mathrm{BA_3}$. Elemental form of $A$ in nature is $\mathrm{A_2}$ (gas) \& that of $B$ is $\mathrm{B_2}$ (solid). Then number of lone pairs in $\mathrm{BA_3}$ is \underline{\hspace{1cm}}.
\qn{10} How many of the following are correct statements.

(i) $B$ is $sp^3$ hybridized in $\mathrm{B_2H_6}$

(ii) $Be$ is $sp^2$ hybridized in $\mathrm{Be_2Cl_4}$

(iii) $\mathrm{Al_2Cl_6}$ is non planar

(iv) Maximum six atoms in a plane in $\mathrm{Al_2Cl_6}$

(v) All Al--Cl length are identical in $\mathrm{Al_2Cl_6}$

(vi) In $\mathrm{[BH_2(NH_3)_2]^+}$ $\mathrm{[BH_4]^-}$ both boron have same hybridization.
\columnbreak
\partsec{3}{MATHEMATICS}
\sect{SECTION-I}
\qn{1} The complete range of $y=\dfrac{4x^2-3x+1}{4x^2+5x+1}$ is $(x\in R)$
\begin{qoptions}
\opt{A}{$(-\infty,-7]\cup\left[\dfrac19,\infty\right)$}
\opt{B}{$(-\infty,-2]\cup[5,\infty)$}
\opt{C}{$\left(-\infty,-\dfrac13\right]\cup[7,\infty)$}
\opt{D}{$(-\infty,-4]\cup[0,\infty)$}
\end{qoptions}
\qn{2} If $\tan A=\dfrac1{\sqrt5}$, then $\cos2A+\tan^2(2A)$ is
\begin{qoptions}
\opt{A}{$\dfrac{23}{12}$}
\opt{B}{$\dfrac1{12}$}
\opt{C}{$\dfrac7{23}$}
\opt{D}{$\dfrac{25}{29}$}
\end{qoptions}
\qn{3} Locus of centroid of triangle formed by $A(2\cos\theta,3\sin\theta)$, $B(-\cos\theta,0)$, $C(0,2\sin\theta)$ is (where $\theta$ is variable)
\begin{qoptions}
\opt{A}{$x^2+9y^2=1$}
\opt{B}{$xy=6$}
\opt{C}{$x^2-25y^2=9$}
\opt{D}{$9x^2+\dfrac{9y^2}{25}=1$}
\end{qoptions}
\qn{4} The value of $\lambda$ for which the point $(\lambda,2)$ is are interior point of $\Delta ABC$ formed by $x+y=4$, $3x-7y=8$ and $4x-y=31$ is
\begin{qoptions}
\opt{A}{$\lambda\in\left(\dfrac{22}3,\dfrac{33}4\right)$}
\opt{B}{$\lambda\in\left(\dfrac{21}3,\dfrac{34}3\right)$}
\opt{C}{$\lambda\in\left(\dfrac{20}3,\dfrac{31}4\right)$}
\opt{D}{$\lambda\in\left(\dfrac{21}4,\dfrac{33}5\right)$}
\end{qoptions}
\qn{5} The system of equation

$2x-y+z=0$

$x-2y+z=0$

$mx-y+2z=0$ has infinite number of non trivial solution for
\begin{qoptions}
\opt{A}{$m=2$}
\opt{B}{$m=-5$}
\opt{C}{$m=5$}
\opt{D}{$m=-2$}
\end{qoptions}
\qn{6} Let the chord of contact with respect to point $(x_1,y_1)$ to the hyperbola $\dfrac{x^2}9-\dfrac{y^2}4=1$, touches the circle $x^2+y^2=25$, then $\dfrac{x_1^2}{81}+\dfrac{y_1^2}{16}$ is
\begin{qoptions}
\opt{A}{$\dfrac1{25}$}
\opt{B}{$\dfrac1{16}$}
\opt{C}{$\dfrac1{32}$}
\opt{D}{$\dfrac1{50}$}
\end{qoptions}
\qn{7} One diagonal of a square is $3x-4y+8=0$ and one vertex is $(-1,1)$ then the area of square is
\begin{qoptions}
\opt{A}{$\dfrac1{50}$ sq. units}
\opt{B}{$\dfrac1{25}$ sq. units}
\opt{C}{$\dfrac3{50}$ sq. units}
\opt{D}{$\dfrac2{25}$ sq. units}
\end{qoptions}
\qn{8} The locus of middle point of chord of parabola $y^2=4x$ which is parallel to the line $y=5x$ is
\begin{qoptions}
\opt{A}{$x^2+y^2=7$}
\opt{B}{$y=\dfrac25$}
\opt{C}{$x=\dfrac32$}
\opt{D}{$x+y=\dfrac56$}
\end{qoptions}
\qn{9} If the line $y=x$ intersect the hyperbola $\dfrac{x^2}9-\dfrac{y^2}{25}=1$ at the points $P$ and $Q$, then the eccentricity of the ellipse with $PQ$ as major axis and minor axis of length $\dfrac5{\sqrt2}$
\begin{qoptions}
\opt{A}{$\dfrac{\sqrt8}3$}
\opt{B}{$\dfrac{\sqrt{10}}3$}
\opt{C}{$\dfrac{\sqrt{12}}3$}
\opt{D}{$\dfrac{\sqrt6}3$}
\end{qoptions}
\qn{10} Eccentricity of an ellipse whose length of latus rectum is one third of length of its minor axis, is
\begin{qoptions}
\opt{A}{$\dfrac12$}
\opt{B}{$\dfrac1{\sqrt3}$}
\opt{C}{$\dfrac{2\sqrt2}3$}
\opt{D}{$\dfrac{5\sqrt5}6$}
\end{qoptions}
\qn{11} Tangent are drawn from the points on the circle $x^2+y^2=41$ to the ellipse $\dfrac{x^2}{16}+\dfrac{y^2}{25}=1$, then the angle between the tangents is
\begin{qoptions}
\opt{A}{$\dfrac\pi4$}
\opt{B}{$\dfrac\pi2$}
\opt{C}{$\dfrac\pi6$}
\opt{D}{$\dfrac\pi{12}$}
\end{qoptions}
\qn{12} If $\sin\theta+\sin\phi=\dfrac14$ and $\cos\theta+\cos\phi=\dfrac15$, then value of $41\sin(\theta+\phi)$ is -
\begin{qoptions}
\opt{A}{40}
\opt{B}{$\dfrac1{16}$}
\opt{C}{11}
\opt{D}{$\dfrac1{50}$}
\end{qoptions}
\qn{13} Minimum value of the sum of the numbers $a^{-5}$, $a^{-4}$, $2a^{-3}$, 6, $a^8$ and $a^7$ is
\begin{qoptions}
\opt{A}{5}
\opt{B}{6}
\opt{C}{11}
\opt{D}{12}
\end{qoptions}
\qn{14} If the circle $x^2+y^2+(4\cos\theta)x+(6\sin\theta)y+\lambda=0$ bisects the circumference of the circle $x^2+y^2-4x-6y+1=0$ (where $\theta,\lambda\in R$), then maximum value of $\lambda$ is
\begin{qoptions}
\opt{A}{$17+\sqrt5$}
\opt{B}{$-\sqrt{50}+48$}
\opt{C}{$\sqrt{98}+5$}
\opt{D}{$2\left(\sqrt{97}\right)-25$}
\end{qoptions}
\qn{15} Let the normal of the parabola $y^2=4ax$ at the point $P(at_1^2,2at_1)$, intersects the curve again at $Q(at_2^2,2at_2)$, then $t_2+t_1+\dfrac2{t_1}$ is
\begin{qoptions}
\opt{A}{$-2$}
\opt{B}{0}
\opt{C}{1}
\opt{D}{$-3$}
\end{qoptions}
\qn{16} Let $x>0$ then minimum value of $\left(x^5+\dfrac{64}{x^5}\right)+10$ is
\begin{qoptions}
\opt{A}{35}
\opt{B}{26}
\opt{C}{75}
\opt{D}{50}
\end{qoptions}
\qn{17} $\dfrac{\tan(\ln15)\tan(\ln3)\tan(\ln5)}{\tan(\ln15)-\tan(\ln3)-\tan(\ln5)}$ is equal to
\begin{qoptions}
\opt{A}{1}
\opt{B}{0}
\opt{C}{$-1$}
\opt{D}{$\dfrac12$}
\end{qoptions}
\qn{18} Number of solutions of the equation $\cos5x=-\cos3x$ in $x\in[0,\pi]$ is
\begin{qoptions}
\opt{A}{3}
\opt{B}{6}
\opt{C}{5}
\opt{D}{7}
\end{qoptions}
\qn{19} If $\log_3x^4=\log_3x^3+\log_37-\log_32$, then $x$ is equal to
\begin{qoptions}
\opt{A}{3.25}
\opt{B}{3.50}
\opt{C}{4.75}
\opt{D}{6.25}
\end{qoptions}
\qn{20} Let two circles intersect each other at $30^\circ$. If $C_1P=1$, $C_2P=2$, then $C_1C_2$ is

(Where $C_1$, $C_2$ are centres of circles and $P$ is their point of intersection)
\begin{qoptions}
\opt{A}{$\sqrt{5-2\sqrt3}$}
\opt{B}{$\sqrt{10+\sqrt5}$}
\opt{C}{$\sqrt{7+\sqrt6}$}
\opt{D}{$\sqrt{5+\sqrt{10}}$}
\end{qoptions}
\sect{SECTION-II}
\qn{1} In a $\Delta ABC$, with usual notations, $4b+4c=5a$, then $9\left(\tan\dfrac C2\tan\dfrac B2\right)$ is
\qn{2} If product of the roots of $\begin{vmatrix}1+x&1&1\\3&1&2\\3&x&3\end{vmatrix}=0$ is $k$, then the value of $2k$ is
\qn{3} For the parabola $y^2=8x$, if length of latus ractum is $\lambda_1$, distance between focus and directrix is $\lambda_2$ and distance between focus and $(6,3)$ is $\lambda_3$, then $(\lambda_1+\lambda_2+\lambda_3)$ is
\qn{4} Let lines through $P(a,2)$, $|a|>3$ meets the ellipse $\dfrac{x^2}9+\dfrac{y^2}4=1$ at $A$ and $D$ and meets the axes at $B$ and $C$ such that $PA\times PD=PC\times PB$. If complete set of $a$ is $(-\infty,-k]\cup[k,\infty)$, then $k$ is
\qn{5} In $\Delta ABC$ with usual notations, let the length of inradius is 2 and sum of coefficient in the expansion of $\left(\left(\dfrac1{r_1}+\dfrac1{r_2}\right)x+\dfrac y{r_3}\right)^{20}$ is $\lambda$ then $50(\lambda\times2^{19})$ is (if expansion is written in powers of $x$ and $y$)
\qn{6} A variable chord of the hyperbola $\dfrac{x^2}4-\dfrac{y^2}8=1$, subtends a right angle at the centre of the hyperbola. If this chord touches a fixed circle concentric with the hyperbola, then the square of radius of the fixed circle is equal to \underline{\hspace{1cm}}
\qn{7} Number of solution(s) of $x$ in $[0,2\pi]$ for the equation $\displaystyle\sum_{r=1}^{10}\sin^2(rx)=10$ is
\qn{8} Let the tangent $y=2x+c$ of hyperbola $\dfrac{x^2}5-\dfrac{y^2}1=1$ intersects $x$ axis at $A$ and $y$ axis at $B$ respectively. If area of $\Delta AOB$ is $\lambda$, then $\left|\dfrac{100\lambda}{19}\right|$ is
\qn{9} Let $S_1$ and $S_2$ are focii of hyperbola $\dfrac{x^2}{\left(\alpha+\frac1\alpha\right)^2}-\dfrac{y^2}5=1$; $\alpha>0$. If $P$ is a variable point on the hyperbola, then minimum value of $|S_1P-S_2P|$ is
\qn{10} Numerical value of harmonic mean of roots of $4x^3-7x^2+3x+1=0$ is
\end{multicols}
\end{document}
